% !TeX root = ../main.tex


\section{Methodology}


\begin{figure*}[!t]
    \centering
    \includegraphics[width=0.85\textwidth]{fig/overview.pdf}
    \caption{Overview}
    \label{fig:overview}
\end{figure*}


We discuss the motivation of our method design as the preliminary. Then we elaborate our method \tool.


\subsection{Preliminary}

% We think the not all elements contributes to the copyright.

% https://opensource.stackexchange.com/questions/13758/gpl-v3-code-used-to-train-some-ai-model-can-i-make-any-use-of-this-model-or-do
% https://news.ycombinator.com/item?id=39589099
% Microsoft announces new Copilot Copyright Commitment for customers
% https://blogs.microsoft.com/on-the-issues/2023/09/07/copilot-copyright-commitment-ai-legal-concerns/

% \textbf{Copyright Definition.} 

% \textbf{Code Copyright.}  Source code – simultaneously embodying \textit{human design intent} and \textit{machine-enforced syntax}. Therefore, we should remove \textit{machine-enforced syntax}.

% \textbf{AST Node.}


Source code is considered a ``literary work'' under copyright law. The ``literary work'' can be viewed as a sequence of tokens arranged to reflect the author's intentions. Existing MIA techniques\todo{CITE} treat all source code tokens equally and compute membership probabilities of a LLM based on a sample of these tokens. Generally, it is recognized that elements such as variable names, data and control flows, and API usage together with the involved variables, embody the author's intent and serve as indicators of source code authorship. However, we observe that some tokens primarily serve a grammatical purpose, functioning to complete the syntax constraint of the programming language.

The inherent structural duality of source code—simultaneously embodying \textit{human design intent} and \textit{machine-enforced syntax}—introduces unique challenges for membership inference in code language models. A large-scale crawl of open-source Python repositories shows that over half of all tokens belong to such grammar patterns (e.g., \texttt{for-in}, \texttt{try-except}). These ubiquitous structures raise two practical obstacles: Syntax-Induced Noise Amplification(language-mandated tokens exhibit near-identical distributions across samples, drowning discriminative signals) and Contextual Ambiguity (inability to distinguish syntax-enforced vs. meaningful occurrences). The root challenge stems from code's \textit{mandatory scaffolding} property: any correct implementation embeds intent within rigid grammatical constraints, creating persistent noise for MIA.


\subsection{Syntax Constraints}
% “语法要求怎么怎么样，要写什么token的代码”

% Operator 常见标点与运算符、属性访问中的点号 Sec. 2.5
% Compound Statement Sec. 8 
% Simple Statement Sec. 7


To address these issues, we construct an \emph{identifier-aware syntax mask} to filter out grammar noise while preserving semantic integrity. The process begins by parsing the code snippet using Python 3.11's \texttt{ast} module to extract all variable, function, and class names, mapping them to their exact character spans in the source via \texttt{lineno} and \texttt{col\_offset}. These spans form a whitelist $\mathcal{W}$ that protects identifiers from removal. Next, we tokenize the code with \texttt{return\_offsets\_mapping=True} to align each subword token with its source offsets. We then initialize a boolean mask $\mathbf{m}$ to true (keep) and set entries to false for tokens matching grammar categories: formatting placeholders (e.g., \texttt{[PAD]}, indents), closing keywords (\texttt{else}, \texttt{except}, etc.), structural symbols (brackets, colons, etc.), punctuation/operators, and attribute dots. If a token's span overlaps $\mathcal{W}$, we revert $\mathbf{m}$ to true.


\begin{table*}[!t]
    \centering
    \footnotesize
    \caption{Syntax Conventions}\label{table:monitor-stats}
    \vspace{-10pt}
    \begin{tabular}{m{2.6cm}m{2.4cm}m{7cm}m{3cm}}
        \toprule
        \textbf{Category} &\textbf{Syntax Node}  & \textbf{Constraint Description} & \textbf{Syntax Token} \\
        \midrule
        Data Model & Lists, Dict, String & Must terminate with appropriate delimiters or accessors for syntactic completeness (e.g., closing brackets or dot notation) & ``]'', ``\}'', ``.'' \\\hline
        Expressions & \texttt{call} expression & Must follow function identifiers and include a parenthesized, comma-separated list of arguments (which may be empty) & ``,'' \\\hline
        \multirow{10}{*}{Single Statements} 
        & \texttt{return} & Probably appear within a function and typically at its end to return a value & return \\
        & \shortstack{\texttt{break} \\\texttt{continue}} & Probably be enclosed within loops (\texttt{for} or \texttt{while}) to alter control flow & break, continue \\
        & \texttt{pass} & Probably be used for completing syntactic structure & pass \\
        & \texttt{raise} & Probably be used within an exception handler to re-raise an exception & raise \\
        & \texttt{import} & Probably with aliases when including external packages  & import, as \\
        & \texttt{with} & Probably with aliases when wrapping a context-manager-defined code  & with, as \\
        \hline
        \multirow{4}{*}{Compound Statements} 
        & \texttt{if} & Must be followed by a colon and indented block; Probably extended with \texttt{elif} or \texttt{else} branches & ``:'', indentation, elif, else, line break \\
        & \texttt{try} & Must be followed by a colon and indented block; Must include at least one of: \texttt{except}, \texttt{except*}, \texttt{finally}, or \texttt{else} & ``:'', indentation, except, except*, finally, else \\
        & \texttt{class} definition & Must follow the class keyword with an identifier and optional parent class list in parentheses, followed by a colon and block & ``)'', ``,'' \\
        & \texttt{method} definition & Must define a function within a class using \texttt{def}, followed by parameters in parentheses and a colon & ``)'', ``,'' \\
        \bottomrule
        \noalign{\hrule height 1pt}
    \end{tabular}
\end{table*}

% \begin{table*}[!t]
%     \centering
%     \footnotesize
%     \caption{Syntax Conventions}\label{table:monitor-stats}
%     \vspace{-10pt}
%     \begin{tabular}{m{2.6cm}m{3cm}m{7cm}m{3cm}}
%         % \noalign{\hrule height 1pt}
%         \toprule
%         \textbf{Category} &\textbf{Syntax Node}  & \textbf{Constraint Description} & \textbf{Syntax Token} \\
%         \midrule
%         Data Model &  Lists, Dict, String & accessing objects, completing array, list, and dict definition & ``]'', ``\}'' , ``.''\\\hline
%         Expressions  &  \texttt{call} expression &  a function call usually accompanied with a possibly empty series of arguments &  ``,''\\     \hline

%         \multirow{5}{*}{Single Statements} & \texttt{return}  & Likely to show up in the end of the function & return \\
%                           & \texttt{break}, \texttt{continue}  & Likely to show up in the end of the \texttt{for}, \texttt{while} loop & break, continue\\ 
%                           & \texttt{pass} & null operation as a placeholder to make syntactically right & pass\\ 
%                           & \texttt{raise}  &  re-raises the exception that is currently being handled in the \texttt{except} clause & raise  \\
%                           & \texttt{import} & import a particular package & import, as \\
%                           & \texttt{with} & wrap the execution of a block with context-manager-defined methods & with as\\
%         \hline
%         \multirow{4}{*}{Compound Statements} & \texttt{if}  & Usually accompanied with \texttt{elif}, \texttt{else} & ``:'', indentation, elif, else, line break, \\
%                             & \texttt{try} &  accompanied with \texttt{except}, \texttt{finally} & ``:'', indentation, except, except*, finally, else \\
%                             & \texttt{class} definition &  syntax convention of defining a class     &  ``)'', ``,'' \\
%                             & \texttt{method} definition &     syntax convention of defining a method   & ``)'', ``,'' \\
%         \bottomrule
%         \noalign{\hrule height 1pt}
%     \end{tabular}
% \end{table*}


                    % +=*/
% A. 
% 7 Single Statements = 再加控制流终结词 `return, break, continue, pass, raise`。

% B. 
% 8 Compound Statements =  % - **闭合关键词**：`else, elif, except, finally` 等；
    % 2.1.8 Indentation = 缩进空格／制表符生成的 token。
    %  `:`, A clause consists of a header and a ‘suite.’ 
% C.
% Data Model
    % 3.1.3. Lists = ]
    % Dictionary = {}, string = '', "", ```'''
% D. Semantic
    % variable convention "_",  variable tokens, - **变量保护**：如果 token 所在区间与步骤 1 得到的任何标识符区间重叠，则**无论上述条件如何**都强制保持 True

% 6.2.3  Parenthesized forms =  )

% E. 2.6 Delimiter
    % , . +=*/
        % - **属性访问中的点号**：`self.foo` 或 `obj.bar` 里，用来连接标识符的 `.` 也被屏蔽。
% ---------------

% - **排版与占位**：`[PAD]`、
% "纯数字"。
                            % - **常见标点与运算符**：逗号,、点号.、'引号、加减乘除等。

% 
\begin{itemize}
    \item \textit{Data Model:}
    \item \textit{Expressions:} 
    \item \textit{Expressions (Sec. ):} parenthesized expression (`()'), 
    \item \textit{Simple Statement (Sec. XX)}: \texttt{return}, \texttt{break}, \texttt{continue}, \texttt{pass}, \texttt{raise}.
    \item \textit{Compound Statement (Sec. XX):} \texttt{else}, \texttt{elif}, \texttt{except}, \texttt{finally}
\end{itemize}


\begin{table}[!t]
    \centering
    \footnotesize
    \caption{Authorship-Reflecting Elements}\label{table:human}
    \vspace{-10pt}
    \begin{tabular}{m{2cm}m{2cm}m{3cm}}
        \toprule
        \textbf{Category} &\textbf{Syntax Node}  & \textbf{Exemplary Token} \\
        \midrule
        Author-defined variables & identifier & hello\_world \\\hline
        Author-defined literals & \texttt{string}, \texttt{bytes}, \texttt{numeric}  & ``OxBEEFXXXX'' \\\hline
    \end{tabular}
\end{table}


\subsection{Identifiable Design Copyright}

String and bytes literals.

There are three types of numeric literals: integers, floating point numbers, and imaginary numbers.

$\mathcal{W}$

Our approach exploits code's dual predictability—grammar tokens are highly deterministic while author choices reveal memorization. By filtering the former, we enhance MIA signals without losing semantic content.


\subsection{Min-k\% Prob with Syntax-Aware Filtering}
We propose \textbf{Syntax-Aware Min-k\% Prob}, an enhanced membership inference metric tailored for code.  After applying the identifier–aware mask, we retain only the log-probabilities $\ell_i=-\log p(t_i\,|\,t_{<i})$ at positions where $\mathbf{m}[i]=1$ and renormalise them to $\tilde{\ell}_i$.  Tokens are ordered by $\tilde{\ell}_i$ and the average of the hardest $k\%$ constitutes the score:
\[
\text{SAMin-}k\%(x)=\frac{1}{|\mathcal{S}|}\sum_{i\in\mathcal{S}}\tilde{\ell}_i, \qquad \mathcal{S}=\text{top-}k\%\big(\{\tilde{\ell}_i\}\big).
\]
Filtering grammar tokens ensures the score is dominated by author–specific choices rather than predictable scaffolding.  The whitelist constraint further guarantees that an identifier, even when split into multiple sub-tokens, is scored as a whole rather than partially suppressed.

\paragraph{Renormalisation.}  Let $\mathcal{K}=\{i\mid \mathbf{m}[i]=1\}$ denote the indices kept by the mask.  For each $i\in\mathcal{K}$ we compute a length-normalised log-probability
\[
\tilde{\ell}_i\;=\;\frac{\ell_i}{|\mathcal{K}|}, \qquad \ell_i=-\log p(t_i|t_{<i}).
\]
This normalisation removes bias from variable sequence lengths once grammar tokens have been excised.

Syntax-aware masking converts a token-level metric into a semantics-level one: only tokens that could plausibly encode authorship survive, yet no identifier is ever broken.  This design preserves fidelity to the original Min-$k\%$ principle while eliminating its chief source of noise in code.
